\documentclass[10pt,a4paper]{amsart}
\usepackage[margin=19mm]{geometry}
\usepackage{amsmath,amssymb,mathtools}
\usepackage[scr=rsfs,cal=euler]{mathalfa}
\usepackage{xfrac}
\usepackage[unicode,hidelinks,bookmarks=false]{hyperref}
\newcommand{\ie}{i.e.\ }
\newcommand{\cf}{cf.\ }
\newcommand{\resp}{respectively\ }
\newcommand{\Subsection}{\subsection}
\providecommand{\nobreakdash}{\nobreak\hbox{-}}
\setlength{\emergencystretch}{2em}
\multlinegap0pt
\usepackage{bm}
\usepackage{bbm}
\usepackage{dsfont}
\usepackage{xspace}
\usepackage{cancel}
\usepackage{mathtools}
\usepackage{amsthm}
\makeatletter
\@ifundefined{lemma}{\newtheorem{lemma}{Lemma}}{}
\renewenvironment{proof}[1]{\par\medskip\noindent\textbf{#1}\enspace\ignorespaces}{\par\medskip}
\makeatother
\makeatletter
\def\ltwonorm#1{\left\lVert #1\right\rVert _{\mu^\star\to\mu^\star}}%
\newcommand{\muhat}{\widehat{\mu}}
\newcommand{\one}{\mathbf{1}}
\expandafter\ifx\csname restatedthm\endcsname\relax
\newenvironment{restatedthm}[2]{%
	\par\addvspace{6pt}%
	\noindent\textbf{#1~\hyperref[#2]{\ref*{#2}}}\ \ \itshape
}{%
	\par\addvspace{6pt}%
}
\fi
\newcommand{\uhat}{\widehat{U}}
\newcommand{\zero}{\mathbf{0}}
\makeatother
\title[Lyapunov-Based Sample Complexity Analysis for Weakly-Coupled]{Lyapunov-Based Sample Complexity Analysis for Weakly-Coupled MDPs}
\author{Tianhao Wu, Matthew Zurek, Weina Wang, Qiaomin Xie}
\date{Source: arXiv:2606.14095v2 (CC BY 4.0)}
\begin{document}
\maketitle

\noindent\emph{Source and licence.} Lyapunov-Based Sample Complexity Analysis for Weakly-Coupled MDPs. Version arXiv:2606.14095v2 (CC BY 4.0), distributed under the Creative Commons Attribution 4.0 International licence, as recorded in the arXiv licence metadata for this exact version and shown on the abstract page https://arxiv.org/abs/2606.14095v2. Full rights notice: licenses/arxiv-2606.14095-CC-BY-4.0.txt.\par\smallskip

\noindent\textit{Editorial scope.} Counted unit: the Lemma whose source label is \texttt{lem:perturbed-Uhat-Huhat-identities} in the article (source0.tex lines 2623--2642, source label \texttt{lem:perturbed-Uhat-Huhat-identities}), delivered together with the article's own complete original proof of it (source0.tex lines 2645--2734, one \texttt{proof} environment of 90 lines). No mathematics is written, repaired or re-derived: statement and proof are the article's own text line for line. Required context retained uncounted: eq:muhat-definition (lines 1550--1553); lem:well-define-Huhat (lines 1573--1584); eq:well-define-Huhat (lines 1576--1578); lem:basic-U-HU-identities (lines 2585--2589). Every definition, display and label of those lines is retained unchanged. Omitted from this package and not redistributed: 1--1549, 1554--1572, 1579--2584, 2590--2622, 2643--2644, 2735--4229. Nothing inside a retained excerpt is omitted or altered: every retained body is delivered exactly as the article's lines are written, so no omission receipt of any kind accompanies this package. Citations: the retained excerpts carry no citation command at all, so no bibliography entry of the article is transcribed and no reference is redistributed. Reference adaptation: none was needed and none was made; every reference inside the delivered unit resolves inside it, so no unresolved reference or citation is produced. Printed slips kept literal: no printed word, symbol, hypothesis or number of the counted statement or of its proof was altered. Layout and class: the article's own class is \texttt{informs3} with packages including theorem, tgtermes, newtxtext, newtxmath, bm, endnotes, inputenc, enumitem, ulem, hyperref, booktabs, xcolor, fontenc, url; the wrapper is the lane's pinned \texttt{amsart} editorial wrapper at 10pt on A4, which restates the theorem-like environments the retained text needs and adds the article's own macro definitions that the retained text needs (\texttt{\textbackslash ltwonorm}, \texttt{\textbackslash muhat}, \texttt{\textbackslash one}, \texttt{\textbackslash uhat}, \texttt{\textbackslash zero}), with extra wrapper packages (bm, bbm, dsfont, xspace, cancel, mathtools). The delivered numbering is the unit's own. Assets: no retained span contains a figure environment, a graphics-inclusion command, a PDF-page inclusion, a table or a TikZ picture; no image file is copied into this package, the package declares no asset dependency and no image file is redistributed with it. Licence: version arXiv:2606.14095v2 of this article is distributed under the Creative Commons Attribution 4.0 International licence, as recorded in the arXiv licence metadata for this exact version and shown on the abstract page https://arxiv.org/abs/2606.14095v2. A full rights notice is delivered at licenses/arxiv-2606.14095-CC-BY-4.0.txt. Characters: every retained excerpt is pure ASCII, so the delivered engine, XeLaTeX with its default Latin Modern text font, reports no missing character.
\begin{equation}
    \muhat\triangleq \mu^\star (I - (\uhat-U)H_U)^{-1}.
    \label{eq:muhat-definition}
\end{equation} 

\begin{lemma}
\label{lem:well-define-Huhat}
    Suppose that $\ltwonorm{U - \one \mu^\star} < 1$ and 
    \begin{equation}
        \ltwonorm{ \uhat-U} \leq \frac{1}{2}\frac{1 - \ltwonorm{U - \one \mu^\star}}{1 + \frac{\ltwonorm{H_{U}}}{\sqrt{\mu^\star_{\min}}} }. \label{eq:well-define-Huhat}
    \end{equation}
    Then
    $
        1-\ltwonorm{\uhat - \one \muhat} \geq \frac{1-\ltwonorm{U - \one \mu^\star}}{2} > 0
    $
    and thus $(I - (\uhat - \one \muhat))$ is invertible.
\end{lemma}

\begin{lemma}
\label{lem:basic-U-HU-identities}
We have \(U\one=\one\) and \(\uhat\one=\one\). Moreover, if
\(I-(U-\one\mu^\star)\) is invertible, then \(H_U\one=\zero\).
\end{lemma}

\begin{lemma}
\label{lem:perturbed-Uhat-Huhat-identities}
Suppose that \(\ltwonorm{U-\one\mu^\star}<1\) and
\eqref{eq:well-define-Huhat} holds. Then
\(
    \muhat\uhat=\muhat,
    \muhat\one=1,
\) and
\(
    H_{\uhat}\one=\zero,
    \muhat H_{\uhat}=\zero,
    (I-\uhat)H_{\uhat}=I-\one\muhat.
\)
Moreover, there exists a row vector \(x\) such that
\(
    H_{\uhat}-H_U
    =
    H_U(\uhat-U)H_{\uhat}+\one x.
\)
\end{lemma}

\begin{proof}{Proof.}
    Abbreviate $X = (I - (\uhat - U)H_U)^{-1}$. By the definition~\eqref{eq:muhat-definition} of $\muhat$, to show $\muhat \uhat = \muhat$ it suffices to show that
        $
            \mu^\star X \uhat = \mu^\star X.
        $
        From the definition of $X$ we have that
        \begin{align*}
            X (I - (\uhat - U)H_U) = I 
            \implies &  X + XUH_U - I = X \uhat H_U \\
            \implies & (X + XUH_U - I)(I-U) = X \uhat H_U(I-U) = X\uhat (I - \one \mu^\star) \\
            \implies & X\uhat = (X + XUH_U - I)(I-U) + X\uhat \one \mu^\star \\
            & \quad\;\;\, = X - XU+XUH_U - XUH_U U - I + U + X\uhat \one \mu^\star,
        \end{align*}
        where we used the property that $H_U(I-U) = I - \one \mu^\star$.
        Now left-multiplying by $\mu^\star$, and also using the facts that $\uhat \one = \one$ and $\uhat\one = U\one = \one$ from Lemma \ref{lem:basic-U-HU-identities}, we have that
        \begin{align*}
            \mu^\star X \uhat 
            &= \mu^\star \left( X - XU+XUH_U - XUH_U U - I + U + X\uhat \one \mu^\star\right) \\
            &= \mu^\star X + \mu^\star XU (-I + H_U -  H_UU + \one \mu^\star) 
            = \mu^\star X,
        \end{align*}
        where in the final equality we used that $-I + H_U -  H_UU + \one \mu^\star = 0$, which is equivalent to the fact that $H_U(I - U) = I - \one \mu^\star$.
        Hence we have that $\muhat \uhat = \muhat$ as desired.

        Now we check that $\muhat \one = \one$. The condition $\ltwonorm{(\uhat-U)H_U}< 1$ (which is implied by~\eqref{eq:well-define-Huhat} as mentioned above) implies the Neumann series expansion $(I - (\uhat-U)H_U)^{-1} = \sum_{t \geq 0} ((\uhat-U)H_U)^t$. Hence by the definition~\eqref{eq:muhat-definition} of $\muhat$ we have
        \begin{align*}
            \muhat \one 
            = \mu^\star (I - (\uhat-U)H_U)^{-1} \one 
            &= \sum_{t \geq 0} \mu^\star ((\uhat-U)H_U)^t \one = \mu^\star \one + \sum_{t \geq 0} \mu^\star ((\uhat-U)H_U)^t (\uhat-U)H_U \one 
            = \one ,
        \end{align*}
        where in the final equality we used that $\mu^\star \one = \one$ and that $H_U \one = \zero$ (from Lemma \ref{lem:basic-U-HU-identities}).

        Having checked $\muhat \one = 1$, we can use an analogous proof as for Lemma \ref{lem:basic-U-HU-identities} to show $H_{\uhat} \one = \zero$.

        Next we check that $\muhat H_{\uhat} = \zero$. From the Neumann series for $(I - (\uhat - \one \muhat))^{-1}$ (since $\ltwonorm{\uhat - \one \muhat}<1$ by Lemma \ref{lem:well-define-Huhat}) we have
        \begin{align*}
            \muhat H_{\uhat} 
            = \muhat (I - (\uhat - \one \muhat))^{-1} - \muhat \one \muhat 
            &= \sum_{t \geq 0} \muhat (\uhat - \one \muhat)^t - \muhat 
            = \muhat I  + \sum_{t \geq 0 }\muhat(\uhat - \one \muhat) (\uhat - \one \muhat)^t - \muhat 
            = \zero,
        \end{align*}
        where we used that $\muhat(\uhat - \one \muhat) = \zero$, which follows from combining the facts $\muhat \uhat = \muhat$ and $ \muhat\one=1$.

        Finally we show that $(I - \uhat)H_{\uhat} = I - \one \muhat$. We have
        \begin{align*}
        (I - \uhat)H_{\uhat} 
        = (I - \uhat + \one \muhat)H_{\uhat} - \one \muhat H_{\uhat}  
        &= (I - \uhat + \one \muhat) \left((I - (\uhat - \one \muhat))^{-1} - \one \muhat \right) - \one \muhat H_{\uhat} \\
        &= I - (I - \uhat + \one \muhat)\one \muhat  - \one \muhat H_{\uhat} 
        = I - \one \muhat+ \one \muhat - \one\muhat ,
    \end{align*}
    where the final equality used $\uhat \one = \one $, $\muhat \one = \one$ and $\muhat H_{\uhat} =\zero$.
    
    Note the resolvent identity $A^{-1} - B^{-1} = A^{-1}(B-A)B^{-1} = B^{-1} (B-A)A^{-1}$ and the fact that $I - (\uhat - \one \muhat)$ and $I - (U - \one \mu^\star)$ are invertible due to the assumptions and Lemma \ref{lem:well-define-Huhat}. We have
\begin{align}
    & (I - (\uhat - \one \muhat))^{-1}  - (I - (U - \one \mu^\star))^{-1} \nonumber \\
    &= (I - (U - \one \mu^\star))^{-1} \left( (I - (U - \one \mu^\star)) - (I - (\uhat - \one \muhat)) \right) (I - (\uhat - \one \muhat))^{-1}\nonumber \\
    &= (I - (U - \one \mu^\star))^{-1} \left( \uhat - U - \one \muhat + \one \mu^\star \right) (I - (\uhat - \one \muhat))^{-1} \nonumber\\
    &= (I - (U - \one \mu^\star))^{-1} \left( \uhat - U \right) (I - (\uhat - \one \muhat))^{-1} - \one (\muhat - \mu^\star) (I - (\uhat - \one \muhat))^{-1} \label{eq:HU-Huhat_step1}
\end{align}
using that $(I - (U - \one \mu^\star))^{-1} \one = \one$ in the final equality, which follows from Lemma \ref{lem:basic-U-HU-identities}.
Next we have 
\begin{multline}
    (I - (U - \one \mu^\star))^{-1} \left( \uhat - U \right) (I - (\uhat - \one \muhat))^{-1} \\
    = H_U  \left( \uhat - U \right) (I - (\uhat - \one \muhat))^{-1} + \one \mu^\star \left( \uhat - U \right) (I - (\uhat - \one \muhat))^{-1} \label{eq:HU-Huhat_step2}
\end{multline}
and
\begin{align}
    H_U  \left( \uhat - U \right) (I - (\uhat - \one \muhat))^{-1} 
    &= H_U  \left( \uhat - U \right) H_{\uhat}  + H_U  \left( \uhat - U \right) \one \muhat \nonumber\\
    &= H_U  \left( \uhat - U \right) H_{\uhat}  + H_U   \zero 
    = H_U  \left( \uhat - U \right) H_{\uhat}, \label{eq:HU-Huhat_step3}
\end{align}
where we used  $U \one = \one $ and $\uhat \one = \one$ from Lemma \ref{lem:basic-U-HU-identities}.
Combining  these calculations, we have that
\begin{align*}
     H_{\uhat} - H_U 
    =& (I - (\uhat - \one \muhat))^{-1} - \one \muhat  - (I - (U - \one \mu^\star))^{-1} + \one \mu^\star  \\
    \stackrel{(i)}{=}& (I - (U - \one \mu^\star))^{-1} \left( \uhat - U \right) (I - (\uhat - \one \muhat))^{-1} - \one (\muhat - \mu^\star) (I - (\uhat - \one \muhat))^{-1}  - \one \muhat + \one \mu^\star  \\
    \stackrel{(ii)}{=}& H_U  \left( \uhat - U \right) (I - (\uhat - \one \muhat))^{-1} + \one \mu^\star \left( \uhat - U \right) (I - (\uhat - \one \muhat))^{-1}  \\
    &- \one (\muhat - \mu^\star) (I - (\uhat - \one \muhat))^{-1}  - \one \muhat + \one \mu^\star \\
    \stackrel{(iii)}{=}& H_U  \left( \uhat - U \right) H_{\uhat} + \one \mu^\star \left( \uhat - U \right) (I - (\uhat - \one \muhat))^{-1}  - \one (\muhat - \mu^\star) (I - (\uhat - \one \muhat))^{-1}  - \one \muhat + \one \mu^\star
\end{align*}
where in steps $(i), (ii),$ and $(iii)$ we used equations~\eqref{eq:HU-Huhat_step1},~\eqref{eq:HU-Huhat_step2}, and~\eqref{eq:HU-Huhat_step3}, respectively. Hence the desired equality is true for
$
    x =  \mu^\star \left( \uhat - U \right) (I - (\uhat - \one \muhat))^{-1}  -  (\muhat - \mu^\star) (I - (\uhat - \one \muhat))^{-1}  -  \muhat +  \mu^\star.
$
\end{proof}

\end{document}
